\section{Standard 4: Consistent Reporting}

\paragraph{Selective Score Reporting}

As evidence for the LLM-As-A-Judge evaluation scores in the original paper's Table 3(b), the first author \href{https://t.me/senior_augur/76}{publicly shared a Telegram link} that showed the higher of two scores was reported for \texttt{min-p} (the reported win rate of $52.01$ corresponds to $p=0.05$, but $p=0.01$ yields a lower win rate of $50.14$) but the lower of two score was reported for \texttt{top-p} (the reported win rate of $50.07$ corresponds to $p=0.9$, but $p=0.98$ yields a higher win rate of $50.43$).

\paragraph{Human Evaluators' Qualitative Responses Fail to Support Claim That \texttt{Min-p} Is Preferred Over Other Samplers}
\label{sec:human_evals:subsec:human_feedback}


At the end of the human evaluation study, the original paper asked human participants to qualitatively describe which sampler(s) they preferred. The paper claimed that human evaluators' qualitative responses support \texttt{min-p} over \texttt{top-p}: ``Participants frequently noted that outputs generated with min-p sampling were more coherent and creative, especially at higher temperatures.''
%
% 
However, we believe that the paper's qualitative human responses contradict this.
We manually annotated the \href{https://github.com/menhguin/minp_paper/blob/main/min_p%20user%20preference%20study%20v2.0%20(Responses)%20-%20Form%20responses%201%20(1)%20(2).csv}{qualitative responses}, then visualized our annotations of the humans' expressed preferences (Fig.~\ref{fig:attentive_human_annotators_preferred_samplers}) and \href{https://docs.google.com/spreadsheets/d/1-KNh4-dbrXBafb9iUQINqFBjDrZvkScPNbGH6IENecQ/edit?gid=1288914057#gid=1288914057}{publicly posted our annotations} in the same format.
We found two results: (1) more human evaluators explicitly preferred \texttt{basic} sampling than preferred \texttt{min-p} sampling, which was not immediately obvious because the \texttt{basic} sampling data were previously excluded (Sec.~\ref{sec:human_evals:subsec:omitted_data}), and (2) \texttt{min-p} was only slightly preferred over \texttt{top-p}.
We provide quotations from human evaluators favoring \texttt{basic} sampling in Appendix~\ref{app:sec:human_qual_responses}.



\subsection{Unsubstantiated and Retracted Adoption Metrics}
\label{sec:community_adoption}

\citet{nguyen2024minp} included specific claims about \texttt{min-p}'s widespread adoption: ``\textbf{Community Adoption:} Min-p sampling has been rapidly adopted by the open-source community, with over 54,000 GitHub repositories using it, amassing a cumulative 1.1 million stars across these projects.''
%
We attempted to verify these numbers through analysis of major GitHub language modeling repositories.
Per our calculations, the combined GitHub stars of leading LM repositories (\texttt{transformers}, \texttt{ollama}, \texttt{llama.cpp}, \texttt{vLLM}, \texttt{Unsloth}, \texttt{mamba}, \texttt{SGLang}, \texttt{llama-cpp-python}) sum to 453k stars as of March 2025, less than half the 1.1M stars claimed by \texttt{min-p} alone.
We could not substantiate either 49k GitHub repositories or 1.1M GitHub stars. 
When we inquired how these numbers were calculated, the authors \href{https://github.com/menhguin/minp_paper/issues/6#issuecomment-2686274162}{publicly stated that GitHub was searched for ``min-p''}, which yields many false positives. The authors retracted both the 54k GitHub repository claim and the 1.1M GitHub stars claim from the ICLR 2025 Camera Ready manuscript.
We highlight this point because 3 of 4 ICLR 2025 reviewers and the Area Chair identified these (retracted) community adoption numbers as the main justification for their strong endorsement.
The ICLR 2025 Camera Ready manuscript has a different statement of community adoption, which we believe remains misleading.
